\documentclass[12pt]{article}
\usepackage[utf8]{inputenc}

\usepackage{_mystyle}

\title{Post RTG Fellowship Essay}
\author{Joseph Sullivan}
\date{October 2025}

\DeclareMathOperator{\Kom}{Kom}
\DeclareMathOperator{\Cyl}{Cyl}
\DeclareMathOperator{\Stab}{Stab}
\DeclareMathOperator{\Slice}{Slice}
\DeclareMathOperator{\chern}{ch}
\DeclareMathOperator{\todd}{td}

\DeclareMathOperator{\ext}{ext}
%\DeclareMathOperator{\hom}{hom}

\newcommand{\cQ}{\mathcal{Q}}
\newcommand{\cZ}{\mathcal{Z}}
\newcommand{\LCom}{\mathcal{LC}}

\newcommand{\rddots}{\text{\reflectbox{$\ddots$}}}


\newcommand{\fg}{\mathfrak{g}}
\newcommand{\fh}{\mathfrak{h}}
\newcommand{\gl}{\operatorname{\mathfrak{gl}}}
\newcommand{\rad}{\mathrm{rad}}

\newcommand{\fa}{\mathfrak{a}}
\newcommand{\fb}{\mathfrak{b}}
\newcommand{\fn}{\mathfrak{n}}
\newcommand{\fs}{\mathfrak{s}}

\DeclareMathOperator{\ad}{ad}
\DeclareMathOperator{\rank}{rank}

\newcommand{\so}{\mathfrak{so}}
\newcommand{\SO}{\mathrm{SO}}
\newcommand{\fsl}{\operatorname{\mathfrak{sl}}}
\DeclareMathOperator{\SL}{SL}
\DeclareMathOperator{\Fl}{Fl}
\DeclareMathOperator{\tr}{tr}

\newcommand{\fW}{\mathfrak{W}}
\newcommand{\cW}{\mathcal{W}}

\DeclareMathOperator{\Gr}{Gr}


\begin{document}

\maketitle

The post-RTG funding would support me to get a solid start on research in my first semester following my oral exam, which I am taking at the end of October. With the extra time aided by the fellowship, I hope to get at a research problem that will eventually become my thesis.

These past semesters I have focused on gaining foundational knowledge in algebraic geometry, as well as pursuing several topics to find my interests. In particular, I have sampled
\begin{itemize}
    \item Syzygies and minimal resolutions
    \item Stability of vector bundles and complexes
    \item Semiorthogonal decompositions of derived categories of sheaves (and related quivers)
    \item The derived category of a quadric, by understanding the spinor bundles through the representation theory of $\so(Q)$.
    \item Some moduli, especially the construction of the Hilbert scheme.
\end{itemize}

Going forwards, I plan to understand the relationship between \textit{pencils} of quadrics and hyperelliptic curves, from a derived category perspective using ``Homological Projective Duality''. This will make good use of my background in derived categories and my focus on quadrics.

Then, there is another angle from which I want to study these intersections of quadrics (and isoptropic Grassmannians) that Aaron has been introducing to me. A theorem of Ramanan-Bhosle describes the intersection of 2 quadrics in $\P^5$ as a certain moduli space of vector bundles on a curve $C$ of genus $2$. Then, this moduli space interacts with certain real/complex representations of the fundamental group of $C$ due to a theorem of Narasimhan-Seshadri.

We are optimistic that we can extend this sort of interaction between intersections of quadrics, moduli of vector bundles, and fundamental group representations to higher dimensions. This project excites me---a lot of machinery that I enjoy is all coming together.

\end{document}
